% RESEARCH_DRAFT_SHA256: 948724c79efbf6fe37ad4d3963c1555dfccf4ae684fa56f9ad29d4855ba3d5b3
% Layout and source-reference wrappers expanded for the web renderer.
\documentclass[11pt]{article}
\usepackage[margin=25mm]{geometry}
\usepackage{amsthm}
\usepackage{fontspec,unicode-math}
\setmainfont{Latin Modern Roman}
\setmathfont{Latin Modern Math}
\usepackage{amsmath,amssymb,mathtools,xurl,hyperref}
\hypersetup{colorlinks=true,urlcolor=blue,linkcolor=blue}
\setcounter{secnumdepth}{0}
\setlength{\parindent}{0pt}
\setlength{\parskip}{5pt}
\setlength{\emergencystretch}{3em}
\newcommand{\sref}[2]{\hyperlink{source:#1:#2}{[S#2]}}
\newcommand{\cataloguescope}[1]{\paragraph{Source record.}#1}
\newcommand{\floor}[1]{\left\lfloor#1\right\rfloor}
\newtheorem{proposition}{Proposition}
\newtheorem{lemma}[proposition]{Lemma}
\begin{document}
\paragraph{Research manuscript.} OpenAI \texttt{gpt-6-astra}, reasoning \texttt{ultra}; 11 October 2026.
\paragraph{Preserved source body.} The complete catalogue mathematical text and sources follow with layout and source links expanded, including source qualifications and its historical preparation notes. New research begins at the section headed Research attempt.

% UnsolvedMath ID 30006102; source catalogue code OWR-14298804-002
\section{Duplicate Strongly Polynomial Linear-Programming Lead-In}\label{30006102}
{\small UnsolvedMath ID 30006102 · Computer Science}
\subsection{Definitions and mathematical statement}

\paragraph{Shared notation.}Standard number systems are $\mathbb N,\mathbb Z,\mathbb Q,\mathbb R,\mathbb C$. Specialized notation and index ranges follow the cited source and the definitions below.


A linear program optimizes $c^Tx$ over a polyhedron $\{x:Ax\leq b\}$. In the simplex method a pivot moves between adjacent basic feasible solutions; a pivot rule chooses the next such move.
\paragraph{Record scope and qualifications.}
The catalogue title is retained as metadata. Named certificates, contracts, bounds or frameworks in that title are not asserted as results of the original source extract.
The supplied registry flags this extraction as unrecoverable. It is displayed as a literal source fragment; the definitions provide context rather than a replacement formal statement.

Source-record limitation: Source flags the statement as unrecoverable: The official report confirms that this extraction is contextual commentary about progress toward the strongly polynomial linear-programming problem represented canonically by OWR-14298804-001. No separate clean statement is promoted; the original record remains publicly visible and retained for provenance.

Related catalogue records: \href{https://mehmetmars7.github.io/Erdosproblems-llm-hunter/problem.html?type=open_problems&id=52}{52}. Separate source IDs are preserved for their provenance.
The source text below is rendered literally where necessary; damaged formulas and composite extracts are not silently repaired into a new mathematical claim.
The supplied catalogue marks this record as unrecoverable. Its extract and terminology are documented, but a complete original problem has not been reconstructed.
\paragraph{Statement recorded in the supplied source.}
This achievement goes beyond previous results, which were limited to feasibility problems, marking a major step toward the aforementioned outstanding open problem. We conclude with mentioning that the workshop included an open problem session, during which six problems were discussed and compiled in this report to guide future research. The report passage is retained verbatim as the relevant context for this extracted problem. \href{https://doi.org/10.4171/owr/2024/50}{[S2]}
\paragraph{Source status.}
The supplied catalogue labels this source record open; that metadata was checked against the supplied snapshot on 2026-10-11. The label is preserved for provenance and is not an independent verification of present open status or mathematical correctness. The catalogue literature assessment is dated 2026-08-22.
\subsection{Short English statement}
The supplied fragment concerns duplicate Strongly Polynomial Linear-Programming Lead-In. Its complete mathematical target is not recovered.
\subsection{Sources}
\begin{itemize}
\item[S1] Ulam AI and the UnsolvedMath Contributors, UnsolvedMath: A Curated Collection of Open Mathematics Problems, record 30006102 (OWR-14298804-002). Catalogue curation and metadata: CC BY 4.0. The source extract is retained; definitions, English summaries and qualifications are editorial additions.
\url{https://huggingface.co/datasets/ulamai/UnsolvedMath}.
\item[S2] Combinatorial Optimization (2025), Oberwolfach Reports, DOI 10.4171/owr/2024/50, pp. 2949--3000
\url{https://doi.org/10.4171/owr/2024/50}.
\item[S3] Dadush et al., Strongly polynomial LP for a restricted class
\url{https://ir.cwi.nl/pub/34290/34290.pdf}.
\item[S4] Oberwolfach Report 50
\url{https://doi.org/10.4171/OWR/2024/50}.
\item Source-specific Oberwolfach Reports licences, including CC BY-SA 4.0 where applicable, are retained; see THIRD\_PARTY\_NOTICES.md.
\end{itemize}
\paragraph{Preparation and review.}
The original source and UnsolvedMath contributors retain attribution for the source extract. Codex (GPT-6; exact serving revision not exposed) assisted with definitions, English summaries, scope notes and formatting. No new human mathematical authorship or independently verified solution is asserted. Permanent source ID 30006102 and code OWR-14298804-002 checked against the complete supplied registry and current repository. Distinct upstream records, including repeated questions and broad or damaged extracts, are retained for complete coverage of the source CS category. A repeated record is not a new independent conjecture. Complete mathematical-equivalence and current-literature review has not been performed. Related canonical identities: 52. Preparation and final review follow contribution prompts C and D, with source gaps disclosed. Structural and rendering checks do not establish a new theorem or open status.
\label{end:30006102}

\section{Research attempt}
\subsection{Literature, source context and scope}
The full catalogue body, including qualifications and historical source status, is preserved above. The original workshop report was inspected in the locally cached publisher PDF text. This attempt distinguishes the source question from any explicitly restricted model introduced below. Repeated records retain their own source IDs and do not become independent conjectures. Elementary results and attributed literature results are not claimed as new original-target progress.

\subsection{Attempted route: repeated feasibility tests recover the optimum}
This is contextual commentary, not a second formal LP conjecture. To investigate its feasibility-versus-optimization distinction, consider an explicitly restricted oracle reduction. An unknown optimum $q$ belongs to
$\{j2^{-B}:0\le j\le2^B\}$. A query at $t$ returns only the bit indicating whether a feasible point with objective at least $t$ exists, equivalently whether $q\ge t$. The reduction must identify $q$ exactly.
\begin{proposition}
Every deterministic comparison tree for this task has worst-case depth at least $\lceil\log_2(2^B+1)\rceil=B+1$.
\end{proposition}
\begin{proof}
There are $2^B+1$ distinct values requiring different outputs. A binary tree of depth at most $r$ has at most $2^r$ leaves, by induction on depth. Two candidate optima ending at the same leaf would give the same output and one would be wrong. Thus $2^r\ge2^B+1$.
\end{proof}
The threshold question can arise from optimizing $x$ over $0\le x\le q$. But when this LP is supplied explicitly, $q$ is already a coefficient and can be read directly. The lower bound concerns the chosen black-box conversion and emphatically is not an LP lower bound. It identifies why a naive threshold-bisection reduction does not inherit a strongly polynomial operation count from a feasibility oracle: the reduction may spend a number of calls depending on precision.
\subsection{Verification and remaining gaps}
At $B=8$ there are 257 possible values and a depth-eight tree has only 256 leaves. Nothing here rules out richer oracle access or a direct strongly polynomial optimization method. The fragment itself has no independent acceptance criterion, and the parent unrestricted LP question is untouched.

\section{References and source attribution}
The primary Oberwolfach workshop report, \url{https://doi.org/10.4171/owr/2024/50}, was inspected on 11 October 2026. Relevant source passages: 2024/50 pp. 2950--2951, contextual lead-in. Catalogue and contributor attribution remain in the original Sources section. Mathematical arguments labeled as this attempt are proved above; cited external results remain attributed to their authors and are not independently reproved or claimed as our discoveries.


\paragraph{Integration review.} Codex (GPT-6; exact serving revision not exposed) reviewed the written arguments and expanded layout and source-reference wrappers. Repository metadata and the production web renderer were checked using Python 3.13 and JavaScript with MathJax 3.2.2. This model-assisted review is not independent human proof verification; no PDF compilation or formal proof checking is asserted.

\section{Completion Estimate}
\noindent\textbf{COMPLETION: 0\%}
No substantiated new progress toward the entire original target is claimed. The worked results and failed route are delimited above; a known result, illustrative model, or auxiliary obstruction is not counted as original-target completion.

\section{Final status and attribution}
\textbf{Attempt status: unresolved.} This describes this research attempt, not a new ruling on the source catalogue's status or on external literature claims.
The mathematics and this self-audit were generated by OpenAI \texttt{gpt-6-astra}, reasoning \texttt{ultra}, on 11 October 2026. Codex (GPT-6; exact serving revision not exposed) orchestrated the batch and contributed formatting/repository review. No human mathematical authorship, independent human proof checking, or formal verification is claimed. The historical preparation paragraphs in the reproduced catalogue body concern that body only.

Contribution prompt A was applied to the complete definition and supplied source record, followed by the adversarial checks of prompt B. Local reproduction files are \texttt{batch011/checks.py}, \texttt{computations.json}, and the per-ID audit JSON. They were not public links at preparation time. Numerical checks are finite and do not prove unrestricted results. No LaTeX engine was found locally; compilation was not performed by this research agent.

\end{document}
